\section{Finite-rank bosonic tomography from complete Husimi shells}
\label{sec:husimi-shell}
The preceding analytic theorem has a direct quantum-information consequence.
The finite parameter is the rank of the density operator, not a photon-number
cutoff or the dimension of the Hilbert space. Exact injectivity, finite-copy
estimation, and stability in trace distance must nevertheless be separated.

\subsection{The infinite-dimensional measurement model}
For $s$ bosonic modes let $\mathcal H=\ell^2(\mathbb N^s)$, with occupation
basis $|\nu\rangle$, and define the normalized coherent vector
\[
 |z\rangle=e^{-\|z\|^2/2}\sum_{\nu\in\mathbb N^s}
               \frac{z^\nu}{\sqrt{\nu!}}|\nu\rangle,
 \qquad z\in\mathbb C^s.
\]
For a positive trace-one operator $\rho$, use the Husimi convention
\[
 Q_\rho(z)=\pi^{-s}\langle z|\rho|z\rangle.
\]
These coherent-state probabilities and their role in tomography are standard;
for a quantitative continuous-variable reconstruction framework see
Landon-Cardinal, Govia and Clerk \cite{LGC2018}. We do not infer a sampling
or numerical reconstruction algorithm from exact identifiability.

For a vector $\psi\in\mathcal H$, the function
\[
 f_\psi(z)=\sum_\nu\overline{\psi_\nu}\frac{z^\nu}{\sqrt{\nu!}}
\]
is entire, since Cauchy--Schwarz bounds its modulus by
$\|\psi\|e^{\|z\|^2/2}$, uniformly on compact sets. If
$\rho=\sum_{j=1}^a\lambda_j|\psi_j\rangle\langle\psi_j|$ has finite rank,
\[
 \pi^s e^{\|z\|^2}Q_\rho(z)
             =\sum_{j=1}^a\lambda_j|f_{\psi_j}(z)|^2.
 \tag{H.1}\label{eq:husimi-bargmann}
\]
No finite mean energy is required for this identity or for the next theorem.

\begin{theorem}[Rank-bounded shell tomography]\label{thm:husimi-shell}
Let $\rho,\sigma$ be bosonic density operators of ranks $a,b<\infty$.
If their Husimi functions have equal radial derivatives of orders
$0,\ldots,m_i-1$ on each of the distinct complete spheres $\|z\|=\rho_i>0$, and
$\sum_i m_i\ge a+b$, then $\rho=\sigma$.
In particular, any $2r$ distinct positive-radius spheres distinguish all
states of rank at most $r$ from one another, without a photon-number cutoff.
\end{theorem}
\begin{proof}
Multiplication by $\pi^s e^{\|z\|^2}$ preserves zero normal multiplicities.
Apply Theorem~\ref{thm:holomorphic-identify} to the maps with components
$\sqrt{\lambda_j}f_{\psi_j}$ and the corresponding components for $\sigma$.
The two polarized kernels coincide. The coefficient of
$z^\nu\overline w^{\mu}/\sqrt{\nu!\mu!}$ is the corresponding density
matrix element, so equality of the entire kernels gives equality of all
matrix elements in the occupation basis.
\end{proof}

The conclusion is rank-$r$ completeness, in the terminology distinguishing it
from rank-$r$ strict completeness \cite{BDK2016}. The latter would distinguish
a rank-$r$ state from every state of arbitrary rank; it is not implied here.
Each shell supplies a full angular function. The theorem does not replace
that continuum by $2r$ scalar observations or claim finite sample complexity.
Indeed no finite collection of coherent-state projectors distinguishes all
pure states in $\mathcal H$: its vector span has an infinite-dimensional
orthogonal complement, containing two orthogonal unit vectors with all those
measurement probabilities zero.

\subsection{Normalized sharp examples and the unknown-rank caveat}
Let $m\ge1$ and define two diagonal, positive, normalized states
\[
 \rho_m^{\mathrm e}=2^{1-m}\sum_{\substack{0\le j\le m\\ j\ {\mathrm{even}}}}
           \binom mj |j\rangle\langle j|,
 \quad
 \rho_m^{\mathrm o}=2^{1-m}\sum_{\substack{0\le j\le m\\ j\ {\mathrm{odd}}}}
           \binom mj |j\rangle\langle j|.
\]
The even and odd binomial sums both equal $2^{m-1}$. Their supports are
orthogonal, and their trace distance is therefore one. Direct evaluation gives
\[
 Q_{\rho_m^{\mathrm e}}(z)-Q_{\rho_m^{\mathrm o}}(z)
     =\frac{2^{1-m}}{\pi}e^{-|z|^2}L_m(|z|^2),
 \qquad L_m(x)=\sum_{j=0}^m(-1)^j\binom mj\frac{x^j}{j!}.
 \tag{H.2}\label{eq:laguerre-ambiguity}
\]
The polynomial $L_m$ has $m$ distinct positive roots. A self-contained argument
uses Rodrigues' identity
$L_m(x)=e^x(d/dx)^m(e^{-x}x^m)/m!$. Integration by parts makes it orthogonal,
with weight $e^{-x}$ on $(0,\infty)$, to every polynomial of degree below $m$.
If it had fewer than $m$ positive sign changes, multiplying by the product of
its sign-changing linear factors would give a nonzero integral of one sign,
contradicting this orthogonality. Its degree is $m$, proving the assertion.

Taking $m=2r-1$ in \eqref{eq:laguerre-ambiguity} gives two rank-$r$ states
indistinguishable on $2r-1$ distinct complete circles, despite trace distance
one. Thus $2r$ is a sharp universal guarantee over radius choices. This
example does not claim failure for every prescribed configuration of $2r-1$
radii. Taking $m=2r$ instead gives ranks $r+1$ and $r$ with identical data on
$2r$ circles. Consequently the universal rank-$r$ completeness guarantee
cannot be relabelled strict completeness when the competitor's rank is unknown.

\subsection{Trace-distance stability needs more than exact shell data}
For states define $D(\rho,\sigma)=\tfrac12\|\rho-\sigma\|_1$. Without an
energy bound no uniform trace-distance modulus exists, even on pure states
and with arbitrarily many fixed bounded measurement circles. The orthogonal
number states $|n\rangle$ and $|n+1\rangle$ have distance one, while on
$|z|\le R$ their Husimi difference is bounded by
\[
 \frac{R^{2n}}{\pi n!}\left(1+\frac{R^2}{n+1}\right)\longrightarrow0.
\]

There is a stronger conditioned example. Fix $E>0$ and take integers $n>2E$.
Set $u=E/n$ and
\[
 |\psi_n^\pm\rangle=\sqrt{1-u}|0\rangle\pm\sqrt u|n\rangle,
 \qquad \rho_n^\pm=|\psi_n^\pm\rangle\langle\psi_n^\pm|.
\]
Both states have rank one, trace one and mean photon number exactly $E$.
Their trace distance and data difference obey
\begin{align}
 D(\rho_n^+,\rho_n^-)&=2\sqrt{u(1-u)},\label{eq:husimi-trace-gap}\\
 \sup_{|z|\le R}|Q_{\rho_n^+}(z)-Q_{\rho_n^-}(z)|
 &\le\frac4\pi\sqrt{u(1-u)}\frac{R^n}{\sqrt{n!}}.
 \label{eq:husimi-data-small}
\end{align}
The first formula follows by diagonalizing the two-by-two difference; the
second follows by expanding the coherent-state amplitudes, whose diagonal
terms cancel. For all sufficiently large $n$, uniformly on the same disk,
\[
 Q_{\rho_n^\pm}(z)\ge\frac{e^{-R^2}}{4\pi}.
\]
Indeed the amplitude before its Gaussian factor is at least
$\sqrt{1-u}-\sqrt u R^n/\sqrt{n!}\ge1/2$.
Thus even a fixed mean-energy bound and a positive lower bound on the measured
intensities do not yield a uniform Holder estimate for trace-distance recovery.
This is a global state error, not the interior intensity error in
Theorem~\ref{thm:shell-instability}.

\begin{proposition}[A compact inverse with a slow necessary modulus]
Fix finitely many complete positive-radius circles, at least $2r$ in number,
and let $\Omega_{r,E}(\delta)$ be the supremum of trace distance over pairs of
rank-at-most-$r$ states with mean photon number at most $E$ and Husimi data
error at most $\delta$ on these circles. Then $\Omega_{r,E}(\delta)\to0$ as
$\delta\downarrow0$. For some $c>0$ and all sufficiently small $\delta$,
\[
 \Omega_{r,E}(\delta)\ge
 c\sqrt{\frac{\log\log(1/\delta)}{\log(1/\delta)}}.
 \tag{H.3}\label{eq:husimi-modulus}
\]
No matching upper rate is asserted.
\end{proposition}
\begin{proof}
The energy constraint makes the state set compact in trace norm. For a direct
verification, let $P_K$ project onto occupations $0,\ldots,K$. Its omitted
probability is at most $E/(K+1)$. By decomposing a state into unit vectors,
using the two-by-two rank-one calculation and Cauchy--Schwarz, one obtains
$\|\rho-P_K\rho P_K\|_1\le2\sqrt{E/(K+1)}$. The finite-dimensional
compressions are totally bounded. Positivity, trace one, rank at most $r$,
and the lower-semicontinuous energy bound are closed conditions, proving
compactness. The Husimi map is continuous into the sup norm because
$|Q_\rho(z)-Q_\sigma(z)|\le\pi^{-1}\|\rho-\sigma\|_1$. It is injective on
this set by Theorem~\ref{thm:husimi-shell}, so its inverse on its image is
uniformly continuous.

For the lower bound put $L=\log(1/\delta)$ and
$n=\lceil4L/\log L\rceil$. The elementary estimate
$\log(n!)\ge n\log n-n+1$ shows that the right side of
\eqref{eq:husimi-data-small} is at most $e^{-L}$ for large $L$ and fixed
$R,E$. Equation~\eqref{eq:husimi-trace-gap} is at least
$\sqrt{2E/n}\ge c\sqrt{\log L/L}$. These rank-one pairs lie in every
rank-at-most-$r$ class, completing the proof.
\end{proof}

All assertions concern ideal complete angular data; no finite-copy confidence
statement follows. The pure-state two-circle case follows from prior scalar
phase retrieval \cite{Perez2021}; bounded-rank and strictly-complete tomography
are established subjects \cite{BDK2016}. The contribution relative to this
project is the no-cutoff finite-rank shell implication, its normalized
Laguerre ambiguities, and the explicit conditioned instability comparison.
Priority for this precise combination has not been established. These
results do not determine the eigenphase statistics of a chaotic propagator.
